% =============================================================================
%  PLAKAT LOEE -- Evolucija.AI 2026
%  Format 70 x 100 cm (pokoncno), ena stran.
%
%  Prevod:   pdflatex plakat-loee.tex        (po potrebi dvakrat)
%            ali:  latexmk -pdf plakat-loee.tex
%
%  Datoteka je samostojna: preambula (mere, barve, pisave, pomozni ukazi) in
%  vsebina sta v isti datoteki. Kazalo:
%     [1] PAKETI IN MERE      -- velikost lista, robovi, sirine stolpcev
%     [2] BARVE               -- celotna paleta
%     [3] VELIKOSTI PISAVE    -- vse velikosti na enem mestu
%     [4] POMOZNI UKAZI       -- \sekcija, \podnapis, \objava ...
%     [5] VSEBINA             -- besedilo in slike plakata
%
%  Slike so v mapi ./slike/ in so vektorske (PDF).
% =============================================================================
\documentclass[final]{article}

% ------------------------------------------------------------ [1] PAKETI IN MERE
\usepackage[T1]{fontenc}
\usepackage[utf8]{inputenc}
\usepackage[slovene]{babel}     % slovensko deljenje besed
\usepackage{graphicx}
\usepackage{xcolor}
\usepackage{tikz}
\usepackage{anyfontsize}        % poljubne velikosti pisave
\usepackage{microtype}
\usepackage{helvet}             % Nimbus Sans (nadomestek za Barlow)
\renewcommand{\familydefault}{\sfdefault}

% Ce imate nalozeno pisavo Barlow (Google Fonts), prevajajte z xelatex ali
% lualatex in odkomentirajte naslednji dve vrstici:
% \usepackage{fontspec}
% \setsansfont{Barlow}[UprightFont=*-Regular, BoldFont=*-Bold, ItalicFont=*-Italic]

% Velikost lista in robovi
\usepackage[paperwidth=70cm, paperheight=100cm,
            left=3.44cm, right=3.44cm, top=2.65cm, bottom=2.22cm]{geometry}

% Sirine stolpcev; vsota levega, presledka in desnega = 63,12 cm (sirina bloka)
\newlength{\levistolpec}\setlength{\levistolpec}{33.20cm}   % besedilo
\newlength{\desnistolpec}\setlength{\desnistolpec}{28.05cm} % slike
\newlength{\presledek}\setlength{\presledek}{1.85cm}
\newlength{\razmikodstavkov}\setlength{\razmikodstavkov}{0.64cm}
\newlength{\razmikslik}\setlength{\razmikslik}{1.10cm}       % razmik med slikami

% ------------------------------------------------------------------ [2] BARVE
\definecolor{ozadje}{HTML}{FCFCFB}       % papir
\definecolor{crnilo}{HTML}{22282C}       % naslov plakata
\definecolor{besedilo}{HTML}{33393D}     % tekoce besedilo
\definecolor{sivo}{HTML}{5F666A}         % podnapisi slik
\definecolor{svetlosivo}{HTML}{6D7479}   % drobni tisk
\definecolor{loeezelena}{HTML}{009933}   % zelena iz logotipa LOEE
\definecolor{temnozelena}{HTML}{00702A}  % zelena za besedilo
\definecolor{ulferdeca}{HTML}{E1352A}    % rdeca UL FE
\definecolor{vmesna}{HTML}{7D8A86}       % sredina barvne crte v glavi
\definecolor{crta}{HTML}{E1E5DF}         % locilne crte

\pagecolor{ozadje}
\pagestyle{empty}
\setlength{\parindent}{0pt}

% -------------------------------------------------------- [3] VELIKOSTI PISAVE
\newcommand{\velNaslov}{\fontsize{90}{97}\selectfont}       % naslov plakata
\newcommand{\velAvtorji}{\fontsize{27}{32}\selectfont}
\newcommand{\velUvod}{\fontsize{30}{42}\selectfont}         % uvodni odstavek
\newcommand{\velBesedilo}{\fontsize{27}{40.5}\selectfont}   % tekoce besedilo
\newcommand{\velSekcija}{\fontsize{24}{28}\selectfont}      % naslov razdelka
\newcommand{\velPodnapis}{\fontsize{19.5}{27.7}\selectfont} % podnapis slike
\newcommand{\velObjava}{\fontsize{18}{24.3}\selectfont}     % seznam objav
\newcommand{\velNoga}{\fontsize{21.75}{28}\selectfont}
\newcommand{\velDrobno}{\fontsize{18}{24}\selectfont}
\newcommand{\velGlava}{\fontsize{19.5}{26}\selectfont}

% --------------------------------------------------------- [4] POMOZNI UKAZI
% Naslov razdelka (zelen, verzalke, spodaj kratka crtica)
\newcommand{\loeesekcija}[2]{%
  \par\vspace{#1}%
  \begingroup\setlength{\parskip}{0pt}%
  {\color{temnozelena}\velSekcija\bfseries\textls[100]{\MakeUppercase{#2}}\par}%
  \vspace{0.24cm}%
  {\color{loeezelena}\rule{2.54cm}{0.106cm}\par}%
  \endgroup
  \vspace{-0.22cm}}
\newcommand{\sekcija}[1]{\loeesekcija{0.94cm}{#1}}      % obicajen razdelek
\newcommand{\prvasekcija}[1]{\loeesekcija{-0.32cm}{#1}} % prvi razdelek v stolpcu

% Podnapis slike: \podnapis{Slika 1}{besedilo}
\newcommand{\podnapis}[2]{\vspace{0.30cm}{\velPodnapis\color{sivo}%
  \textcolor{temnozelena}{\textbf{#1}} $\cdot$ #2\par}}

% Ena bibliografska enota
\newcommand{\objava}[1]{{\velObjava\color{besedilo}#1\par}\vspace{0.26cm}}

% Tanka locilna crta cez celotno sirino
\newcommand{\locnica}{{\color{crta}\rule{\textwidth}{0.053cm}}\par}

% =============================================================================
% [5] VSEBINA
% =============================================================================
\begin{document}

% ---------------------------------------------------------------- GLAVA
\noindent
\begin{minipage}[b]{0.34\textwidth}
  {\color{svetlosivo}\velGlava Katedra za osnove elektrotehnike,\\
   matematiko in fiziko\par}
\end{minipage}\hfill
\begin{minipage}[b]{0.30\textwidth}
  \centering\includegraphics[height=2.54cm]{slike/ulfe-logo.pdf}
\end{minipage}\hfill
\begin{minipage}[b]{0.30\textwidth}
  \raggedleft\includegraphics[height=3.44cm]{slike/loee-logo.pdf}
\end{minipage}

\vspace{0.90cm}
\noindent\begin{tikzpicture}
  \shade[left color=loeezelena, middle color=vmesna, right color=ulferdeca]
        (0,0) rectangle (\textwidth,0.16cm);
\end{tikzpicture}

% ---------------------------------------------------------------- NASLOV
\vspace{1.16cm}
{\color{crnilo}\velNaslov Optimizacija in umetna inteligenca v energetiki\par}

\vspace{0.64cm}
{\color{temnozelena}\velAvtorji\bfseries
 asist.\ Andraž Sivec $\cdot$ asist.\ mag.\ Miran Meža\par}

\vspace{0.79cm}
{\color{sivo}\velUvod Kako naj se baterijski hranilnik obnaša, da bo račun za
elektriko čim nižji -- in koliko od teoretičnega optimuma je sploh dosegljivo,
ko prihodnosti ne poznamo?\par}

\vspace{1.22cm}

% ---------------------------------------------------- LEVI STOLPEC (besedilo)
\noindent
\begin{minipage}[t]{\levistolpec}
  \velBesedilo\color{besedilo}\setlength{\parskip}{\razmikodstavkov}

  \prvasekcija{Laboratorij in industrija}

  Laboratorij za osnove elektrotehnike in elektromagnetiko (LOEE) na Fakulteti
  za elektrotehniko Univerze v Ljubljani povezuje raziskovalno, projektno in
  pedagoško delo; vodi ga prof.\ ddr.\ Iztok Humar.

  Podjetjem in drugim naročnikom so na voljo analize, meritve in zasnova
  rešitev: presoja elektromagnetne kompatibilnosti, vrednotenje elektromagnetnih
  polj v okolici naprav in postrojev, zastiranje polj, indukcijsko segrevanje
  ter razvoj senzorike, vgrajenih sistemov in zajema podatkov. Pot je običajno
  enaka -- meritev, model, prototip in preizkus na realnih podatkih.
  % Tu je smiselno dodati konkretne partnerje ali projekte, npr.:
  % Med naročniki so \textbf{ime podjetja}, \textbf{ime podjetja} in drugi.

  Enako velja za umetno inteligenco: zanima nas, kaj od modela ostane v
  obratovanju, zato ob točnosti merimo tudi ceno -- koliko podatkov, računanja
  in vzdrževanja zahteva ter koliko obljubljene koristi res prinese. Metode
  strojnega učenja uporabljamo pri razvrščanju anomalij v brezžičnih omrežjih,
  pri energijsko varčnem globokem učenju, pri razpoznavanju porabnikov iz
  merilnih podatkov in pri analizi polnjenja električnih vozil -- del tega z
  Institutom »Jožef Stefan«.

  \sekcija{Optimizacija baterijskih hranilnikov}

  Delovanje domačega hranilnika zapišemo kot mešano celoštevilsko linearno
  optimizacijo v 15-minutnih intervalih čez celo leto in ga poženemo na realnih
  meritvah porabe ter na dejanskih ponudbah slovenskega trga -- z omrežnino po
  časovnih blokih, prispevki in dajatvami vred.

  Odgovor ni le »koliko prihranim«, temveč tudi, kje je meja: katera kapaciteta
  se še izplača, koliko je vredna vsaka dodatna kilovatna ura in koliko
  življenjske dobe baterije pri tem porabimo. Razlike med ceniki ponudnikov in
  med tipi gospodinjstev so velike -- največ pridobijo tista s toplotno črpalko,
  električnim avtomobilom in sončno elektrarno.

  \sekcija{Krmiljenje in spodbujevalno učenje}

  Optimum s popolnim poznavanjem prihodnosti je zgornja meja, ki je v resnici ni
  mogoče doseči -- pravi krmilnik ne pozna ne jutrišnjega sonca ne jutrišnje
  porabe. Zato primerjamo preprosta pravila, napovedne modele, optimizacijo na
  drsečem oknu in spodbujevalno učenje (SARSA, DQN) ter merimo, koliko od
  teoretičnega optimuma dejansko ostane.

  \sekcija{Energetske skupnosti}

  Isti model razširimo na skupnosti, kjer si člani delijo proizvodnjo.
  Preverjamo, kako pogodbeni delež deljene energije, izbira sheme in skupni
  hranilnik vplivajo na račun posameznega člana in skupnosti kot celote.

\end{minipage}%
\hspace{\presledek}%
% ---------------------------------------------------- DESNI STOLPEC (slike)
\begin{minipage}[t]{\desnistolpec}
  \color{besedilo}

  \includegraphics[width=\linewidth]{slike/marginal-value.pdf}
  \podnapis{Slika 1}{Koliko je vredna vsaka dodatna kilovatna ura hranilnika pri
  posameznem ceniku. Kjer krivulja pade pod ceno vgrajene kapacitete (vodoravne
  črte), povečevanje baterije nima več ekonomskega smisla.}

  \vspace{\razmikslik}
  \includegraphics[width=\linewidth]{slike/net-benefit.pdf}
  \podnapis{Slika 2}{Neto letna korist glede na velikost hranilnika pri različnih
  cenah vgrajene kapacitete. Optimum je vrh krivulje: dražja baterija pomeni
  manjšo optimalno kapaciteto.}

  \vspace{\razmikslik}
  \includegraphics[width=\linewidth]{slike/share-sweep.pdf}
  \podnapis{Slika 3}{Letni račun energetske skupnosti glede na pogodbeni delež
  deljene energije, z baterijo in brez nje. Označene točke so optimum posamezne
  sheme delitve.}

  % --- Rezervni sliki (odkomentirajte, ce zelite katero namesto zgornjih) ----
  % \vspace{\razmikslik}
  % \includegraphics[width=\linewidth]{slike/achievable-saving.pdf}
  % \podnapis{Slika 4}{Koliko je pri posameznem ceniku sploh mogoče prihraniti z
  % optimalno vodenim hranilnikom v primerjavi z istim gospodinjstvom brez
  % baterije.}
  %
  % Vecpanelna slika poteka dneva; trim = {levo spodaj desno zgoraj} v pt:
  % \includegraphics[trim={0 345.1pt 0 43.3pt}, clip, width=\linewidth]{slike/trace-day.pdf}
  % \includegraphics[trim={0 223.4pt 0 149.0pt}, clip, width=\linewidth]{slike/trace-day.pdf}
  % \includegraphics[trim={0 93.8pt 0 269.0pt}, clip, width=\linewidth]{slike/trace-day.pdf}

\end{minipage}

\vfill   % preostali prostor porine spodnji pas na dno lista

% ---------------------------------------------------------------- OBJAVE
\locnica
\vspace{0.635cm}

\noindent
{\color{temnozelena}\velSekcija\bfseries\textls[100]{IZBRANE OBJAVE}\par}
\vspace{0.30cm}
{\color{loeezelena}\rule{2.54cm}{0.106cm}\par}
\vspace{0.28cm}
{\color{svetlosivo}\velObjava revije s faktorjem vpliva najprej, nato
 konference in domače objave\par}
\vspace{0.42cm}

\objava{B.~Bertalanič, M.~Meža, C.~Fortuna: \textit{Resource-aware time series
  imaging classification for wireless link layer anomalies.} IEEE Transactions
  on Neural Networks and Learning Systems 34(10), 2022.}
\objava{M.~Meža, G.~Strle, M.~Meža: \textit{Survival analysis of electric
  vehicle charging behavior and the temporal evolution of feature effects.}
  Scientific Reports 15, 2025.}
\objava{A.~Pirnat, B.~Bertalanič, G.~Cerar, M.~Mohorčič, M.~Meža, C.~Fortuna:
  \textit{Towards sustainable deep learning for wireless fingerprinting
  localization.} IEEE ICC 2022.}
\objava{A.~Sivec, M.~Meža: \textit{SARSA ali DQN za krmiljenje baterijskega
  hranilnika.} Elektrotehniški vestnik 93(1--2), 2026.}
\objava{A.~Sivec, M.~Meža: \textit{Optimizacija cene električne energije
  gospodinjstva s spodbujevalnim učenjem.} ERK 2025.}
\objava{A.~Sivec, M.~Pantoš, M.~Meža: \textit{Delovanje baterijskih hranilnikov
  na slovenskem energetskem trgu.} V pripravi -- vir slik 1--3.}

% ---------------------------------------------------------------- NOGA
\vspace{0.30cm}
\locnica
\vspace{0.635cm}

\noindent
\begin{minipage}[b]{0.55\textwidth}
  {\velNoga\color{besedilo}%
   \textbf{asist.\ Andraž Sivec} $\cdot$ andraz.sivec@fe.uni-lj.si\\[0.21cm]
   \textbf{asist.\ mag.\ Miran Meža} $\cdot$ miran.meza@fe.uni-lj.si\par}
\end{minipage}\hfill
\begin{minipage}[b]{0.42\textwidth}
  \raggedleft
  {\velDrobno\color{svetlosivo}Predstavljeno na konferenci Evolucija.AI 2026\par}
  \vspace{0.42cm}
  \includegraphics[height=2.75cm]{slike/kcui-logo.pdf}
\end{minipage}

\end{document}
